\section{Tests and Evaluation}
\label{sec:evaluation}

\subsection{Testing Strategy}
\label{sec:strategy}

The code generation is tested at three levels, and every level runs in the
build on every change.

\begin{enumerate}
  \item \textbf{The oracle comparison.} The build runs the whole pipeline on the
    union of the \texttt{minimal} and \texttt{data} projects and compares the
    tree it writes with the committed trees, file for file and byte for byte.
    The comparison sits in a separate test module that holds no EMF type and
    calls no Java code: it reads only the files the pipeline left behind, the
    same interface the production implementation's own test reads. A file the
    run should have written and did not fails as a missing file.
  \item \textbf{The pipeline's entry point.} JUnit tests call
    \texttt{Pipeline.render} for each project's share and for the shared
    files, render the same inputs twice into two directories and compare them,
    and check that a set the constraints refuse writes no file at all.
  \item \textbf{The runner.} The render bundle's own tests run the module on a
    resolved model built in code, check the one file it writes, check that the
    caller's model is left as it was, and run a deliberately broken module to
    check that a parse error fails the run with Acceleo's message.
\end{enumerate}

The Java code of the render bundle is held to 100\% line and branch coverage
and a 100\% mutation score with PIT, like every module of the implementation.
Mutation testing does not reach the templates themselves, so their evidence is
the oracle comparison: a template that writes one character differently fails
it.

\subsection{Results}
\label{sec:results}

Table~\ref{tab:results} shows the result. Every file of every tree the
production implementation generates is equal, and the second render of the
same inputs is byte-identical to the first.

\begin{table}[!htbp]
\centering
\small
\begin{tabular}{>{\raggedright\arraybackslash}p{0.34\linewidth}r>{\raggedright\arraybackslash}p{0.42\linewidth}}
\hline
\textbf{Committed tree} & \textbf{Files} & \textbf{Result} \\
\hline
\path{minimal/rendered/} & 9 & equal byte for byte \\
\path{data/rendered/} & 28 & equal byte for byte \\
\path{_estate/rendered/} & 7 & equal byte for byte \\
\textbf{Total} & \textbf{44} & every file equal; no extra file written \\
\hline
\end{tabular}
\caption{The Rendered Trees Compared With the Oracles}
\label{tab:results}
\end{table}

The 44 files exercise every one of the sixteen kinds of file: a blue-green
workload with its Canary and pod monitor in \texttt{minimal}, and in
\texttt{data} three stateful services with their services, volume claims,
backup jobs, a configuration file, a sidecar container, environment variables
read from synchronised secrets, the secret operator's connection and objects,
and the secret store's policies and roles.

\subsubsection{Loading the Generated Files}
\label{sec:loading}

The task asks for the generated code to be loaded in a tool that handles it.
For Kubernetes files there are two such tools short of a live cluster, and the
tree passes both.

\begin{itemize}
  \item \textbf{kustomize}, which Flux uses to apply a directory, builds every
    directory from its index file: the \texttt{notes} project to 5 objects, the
    \texttt{data} project to 27, and the shared directory of the public entry
    point to 2. A wrong index entry, a missing file or a file that is not
    valid YAML fails the build.
  \item \textbf{kubeconform} validates every object against the schemas of
    Kubernetes~1.31 and, for the custom resources of Flagger, the Prometheus
    operator, Traefik and the Vault secrets operator, against the schemas in
    the public CRD catalogue: all 47 objects in the 40 YAML files are valid,
    with unknown fields refused (\texttt{-strict}).
\end{itemize}

The same kubeconform check is a gate of the repository, run on the committed
trees, so the oracles themselves are known to be valid. Because the generated
tree equals them byte for byte, it inherits that validity; the run above
checks the generated tree directly as well.

\subsubsection{Defects Found by Testing}
\label{sec:defects}

The oracle comparison found every defect below. None of them would have been
visible by reading the generated files, because each produced files that
looked plausible.

\begin{itemize}
  \item \textbf{No file at all.} The first run generated nothing and reported
    no error. The main template's parameter type resolves only against the
    packages the query environment knows, and the target metamodel had not
    been registered with it, so no model element matched the main template.
    The runner now registers the package, and the tests check the files, not
    only the diagnostic.
  \item \textbf{The margin rule.} The first version of the module was written
    as a plain text template and failed to parse, because every body line
    lacked its two spaces of margin (Section~\ref{sec:margins}).
  \item \textbf{A dash on every line.} A container called after \texttt{- }
    turned each of its keys into a separate list item, and so did a route,
    because the text before a call is repeated on every line it writes.
  \item \textbf{Nested indentation.} The environment list, the backup's
    service account and the synchronised secrets came out two columns too
    deep: a block nested in a \texttt{for} and an \texttt{if} has two levels of
    margin, not one.
  \item \textbf{A resource is not a model element.} The first shared-index
    query navigated the resource holding the models, which AQL cannot do; it
    now uses \texttt{allInstances} (Section~\ref{sec:sharedindex}).
  \item \textbf{Two defects in the transformation.} Comparing the
    \texttt{data} tree showed that the transformation chose a stateful
    controller for a process with a volume, where the production implementation and the
    oracle use a deployment, and comparing the Canaries showed that the
    revision differed from the production implementation's
    (Section~\ref{sec:refinement}). Both were fixed in the transformation, not
    in the templates.
\end{itemize}

\subsection{Limitations}
\label{sec:limitations}

\begin{itemize}
  \item \textbf{Two kinds of file are not generated.} A disruption budget and
    a database migration job appear only in the hand-written goal-state trees
    of \texttt{auth} and \texttt{knowledge}. Neither implementation generates
    them yet, and the target metamodel does not describe them.
  \item \textbf{The delivery machinery is not generated.} The processes that
    run the delivery itself (the proxies, Flagger and the release gate) roll
    out in a third way the specification defines, which neither implementation
    renders yet, so the projects holding them have no rendered tree.
  \item \textbf{The quoting rule is partial.} The production serializer also
    quotes a string that starts with an indicator character or holds a colon
    followed by a space. No value in the examples needs that, so the templates
    do not spell it; such a value would be written unquoted and would differ
    from the production file.
  \item \textbf{Some template branches are not exercised by an example.} No
    example has a route that matches an exact path, a duration given in
    minutes or hours, a configuration file without a final line break, or a
    boolean-looking environment value. These branches are written to the same
    rules as the exercised ones, but no oracle proves them.
\end{itemize}

\subsection{Reflection}
\label{sec:reflection}

Comparing byte for byte made the task harder and the result more convincing.
Most of the work was not in choosing which objects to write, which the target
metamodel already decides, but in matching the production serializer's text
exactly: its key order, its quoting and its indentation. That precision is
also what made the comparison find real defects in the transformation, which
a looser comparison would have hidden.

Keeping the templates free of decisions worked as intended. Every time a
template seemed to need a decision, such as which monitor to write or which
identities hold a secret, the decision moved into the transformation and the
target metamodel gained the value, so the templates stayed a pure spelling of
the model. The cost was a larger target metamodel than Task~1 planned.

Acceleo~4's margin rule was the steepest part to learn, because it differs from
every other template language: indentation in the module is both syntax and
output. Once understood, it was also what made the fragments work, because one
template can be placed at any depth by its caller.
